\documentclass{article}
\usepackage{graphicx} % Required for inserting images
\usepackage{glossaries}

\title{Writing Sketches}
\author{Gabriel Luiz Espindola Pedro}
\date{April 2026}

\begin{document}

\maketitle

% Motivação
Comparar municipios é uma tarefa dificil dada as diferentes frentes de análise que podem definir qual município é melhor que outro e porque.  Saber quais municipios se destacam trazem foco à quais devem ser os pontos de atenção para agências reguladoras. Cidades bem qualificadas podem levantar enfoques para melhorias, enquanto cidades mal qualificadas podem deficiências críticas, indicando quais devem ser os passos para evolução de qualidade de serviços de telecomunicações.

% Objetivo
Queremos clasificar os municípios, definindo quais possuem as melhores e as piores qualidade de serviços de telecomunicações. Para isso realizar uma agregação simples a partir de uma combinação linear com dados de métricas QoS deixa de considerar outros fatores que implicam na comparação de qualidade telecomunicações para cada município.

Nosso objetivo é definir um escore padrão para cada município, de modo que a partir dele possamos estabelecer um ranking geral de municípios. Estabelecendo um escore para cada entidade definida por dados de  município por mês podemos analisar a evolução temporal do município, assim como sua competição com outros municípios.


% Manifold
Devemos considerar que, entidades com características similares devem possuir escores próximos, ou seja, estão próximas dentro da geometria espacial onde os dados estão definidos. Em outras palavras, os dados de características das entidades constroem um manifold geométrico, com proximidade local indicando pequena variação de escore.

% Grafo
Este manifold pode ser representado por um grafo construido a partir da proximidade das entidades, onde a estrutura local é representada a partir a conexão entre as entidades, descrita por uma matriz de adjacencia $A$.

% Subsampling
Grandes grafos podem trazer ostensividade computacional, para isso existem técnicas de subsampling analisando um subconjunto dos dados, realizando a amostragem e processamento iterativamente para poder processar o grafo completo.

% Pareto


% Monotonicidade
Visto que maiores 
havendo uma relação de evolução objetiva, ou seja, valores maiores de QoS devem indicar maior escore.

\end{document}